% TODO make hw stylesheet
\documentclass[10pt]{scrartcl}
\usepackage[a4paper, total={6in, 8in}]{geometry}
\usepackage[usenames,dvipsnames,svgnames]{xcolor}
\usepackage[shortlabels]{enumitem}
\usepackage[framemethod=TikZ]{mdframed}
\usepackage{amsmath,amssymb,amsthm}
\usepackage[colorlinks]{hyperref}
\usepackage{multicol}
\usepackage{tabularx}

\usepackage{listings}

\hypersetup{
	colorlinks=true,
	linkcolor=blue,
	filecolor=magenta,      
	urlcolor=magenta,
	pdftitle={CSE 214 HW}
}

\usepackage{microtype}
\usepackage{mathtools}
\usepackage[headsepline]{scrlayer-scrpage}
\usepackage{thmtools}
\usepackage[textsize=footnotesize]{todonotes}


\declaretheoremstyle[spaceabove=6pt,spacebelow=6pt]{basehead}

\declaretheorem[style=basehead,name=Problem]{problem}
\declaretheorem[
	style=basehead,
	thmbox={style=M,bodystyle=\normalfont,headstyle=\normalfont \textbf{Algorithm \thealgo}},
	name=Algorithm,
	within=section
] {algo}
\renewcommand{\thealgo}{\thesection(\roman{algo})}

\newenvironment{alg}[1]{\begin{algo}[#1]
		\setlist{nolistsep}
		\begin{enumerate}[\footnotesize{\arabic*.}\enspace, noitemsep]}{\end{enumerate}\end{algo}}

\addtolength{\textheight}{3.14cm}
\providecommand{\ol}{\overline}
\providecommand{\ul}{\underline}
\providecommand{\eps}{\varepsilon}
\providecommand{\half}{\frac{1}{2}}
\providecommand{\dang}{\measuredangle} %% Directed angle
\providecommand{\CC}{\mathbb C}
\providecommand{\FF}{\mathbb F}
\providecommand{\NN}{\mathbb N}
\providecommand{\QQ}{\mathbb Q}
\providecommand{\RR}{\mathbb R}
\providecommand{\ZZ}{\mathbb Z}
\providecommand{\dg}{^\circ}
\providecommand{\ii}{\item}
\providecommand{\setdiff}{\smallsetminus}
\providecommand{\alert}[1]{{\sffamily\textbf{\textcolor{blue}{#1}}}}
\providecommand{\inv}{^{-1}}
\providecommand{\mb}[1]{\mathbf{#1}}
\providecommand{\dif}{\;d}

\reversemarginpar

\newenvironment{soln}{\begin{proof}[Solution]}{\end{proof}}
\newenvironment{walkthrough}{\noindent \textbf{Walkthrough.}}{\medskip}
\newlist{walk}{enumerate}{3}
\setlist[walk]{label=\bfseries (\alph*)}

\providecommand{\pow}{\operatorname{Pow}}
\providecommand{\vocab}[1]{{\textbf{\textcolor{ForestGreen}{#1}}}}
\providecommand{\lcm}{\operatorname{lcm}}
\providecommand{\abs}[1]{\left|#1\right|}

\usepackage{asymptote}
\begin{asydef}
	defaultpen(fontsize(10pt));
	unitsize(1cm);
	size(8cm); // set a reasonable default
	usepackage("amsmath");
	usepackage("amssymb");
	settings.tex="pdflatex";
	settings.outformat="pdf";
	// Replacement for olympiad+cse5 which is not standard
	import geometry;
	// recalibrate fill and filldraw for conics
	void filldraw(picture pic = currentpicture, conic g, pen fillpen=defaultpen, pen drawpen=defaultpen)
	{ filldraw(pic, (path) g, fillpen, drawpen); }
	void fill(picture pic = currentpicture, conic g, pen p=defaultpen)
	{ filldraw(pic, (path) g, p); }
	// some geometry
	pair foot(pair P, pair A, pair B) { return foot(triangle(A,B,P).VC); }
	pair orthocenter(pair A, pair B, pair C) { return orthocentercenter(A,B,C); }
	pair centroid(pair A, pair B, pair C) { return (A+B+C)/3; }
	// cse5 abbreviations
	path CP(pair P, pair A) { return circle(P, abs(A-P)); }
	path CR(pair P, real r) { return circle(P, r); }
	pair IP(path p, path q) { return intersectionpoints(p,q)[0]; }
	pair OP(path p, path q) { return intersectionpoints(p,q)[1]; }
	path Line(pair A, pair B, real a=0.6, real b=a) { return (a*(A-B)+A)--(b*(B-A)+B); }
	// cse5 more useful functions
	picture CC() {
		picture p=rotate(0)*currentpicture;
		currentpicture.erase();
		return p;
	}
	pair MP(Label s, pair A, pair B = plain.S, pen p = defaultpen) {
		Label L = s;
		L.s = "$"+s.s+"$";
		label(L, A, B, p);
		return A;
	}
	pair Drawing(Label s = "", pair A, pair B = plain.S, pen p = defaultpen) {
		dot(MP(s, A, B, p), p);
		return A;
	}
	path Drawing(path g, pen p = defaultpen, arrowbar ar = None) {
		draw(g, p, ar);
		return g;
	}
\end{asydef}

\usepackage{mathrsfs}

\ihead{\footnotesize CSE 214 HW02}
\ohead{\footnotesize Last compiled \today}

\title{\vspace{-2em}CSE 214 HW02}
\author{\large Aditya Pahuja}
\date{\vspace{-2em}}
\begin{document}
\maketitle

\section{Problem 1}

If an element $x$ in $A$ has the property that $-(2x + 17)$ is also in $A$,
then say that $x$ is \vocab{valid}.

\begin{alg}{\textsc{FindX-Naive($A[1\dots n]$)}}
	\ii valuefound $\gets$ false
	\ii for $i\gets 1$ to $n$ do
	\ii \qquad for $j\gets 1$ to $n$ do
	\ii \qquad\qquad if $A[j] = -(2\cdot A[i] + 17)$ then
	\ii \qquad\qquad\qquad print($A[i]$)
	\ii \qquad\qquad\qquad valuefound $\gets$ true
	\ii \qquad\qquad\qquad break
	\ii if not valuefound then
	\ii \qquad print(``no valid value found'')
\end{alg}

This algorithm is a brute force algorithm: for each $x := A[i]$,
it checks each element of $A$ to see if it is equal to $(-2x + 17)$.
It then prints out the first such $A[i]$ and terminates,
or, if no such $A[i]$ is found, prints that no valid value was found.

This is thus an $\mathcal O(n^2)$ algorithm, since in the worst case,
if our desired value doesn't exist, the algorithm iterates through
all pairs $(i, j)$ with $1\le i,j\le n$ (and processing such a pair
takes $\Theta (1)$ time). Also, the extra space is $\Theta(1)$
because the only extra space used is for the valuefound variable.

\begin{alg}{\textsc{FindX-Better($A[1\dots n]$)}}
	\ii valuefound $\gets$ false
	\ii create array $B[1\dots n]$
	\ii for $i\gets 1$ to $n$ do
	\ii \qquad $B[n+1-i]\gets -(2\cdot A[i] + 17)$
	\ii for $i\gets 1$ to $n$ do
	\ii \qquad $x\gets A[i]$
	\ii \qquad if $\textsc{BinarySearch($B[1\dots n]$, $-(2x + 17)$)}
	> 0$ then
	\ii \qquad \qquad print($x$)
	\ii \qquad \qquad valuefound $\gets$ true
	\ii \qquad \qquad break
	\ii if not valuefound then
	\ii \qquad print(``no valid value found'')
\end{alg}

The way that $B$ is defined is that the last ($n$th) element of $B$ is
$-(2\cdot A[1] + 17)$, the second-to-last ($(n-1)$th) element
is $-(2\cdot A[2] + 17)$ and so on up to the first element $-(2\cdot A[n] + 17)$
--- in particular, $B$ contains $-(2x + 17)$ for every $x$ in $A$.
Since $A$ is sorted, the sequence
$-(2\cdot A[1]+17)$, $-(2\cdot A[2]+17)$, \dots, $-(2\cdot A[n]+17)$
is decreasing ($x\mapsto -(2x+17)$ is clearly decreasing in $x$),
so $B$, whose elements are this sequence in reverse order,
must be a sorted array.

What the algorithm does then is check whether any of the elements of $A$
are also in $B$ (via binary search on each element of $A$,
since $B$ is sorted!).
This is equivalent to the existence of a valid value existing as,
if some number $y$ were in both $A$ and $B$,
then $y$ being in $B$ by definition means that $y = -(2x + 17)$
for some $x$ in $A$, so $x$ and $-(2x + 17)$ are both in $A$,
hence $x$ is a valid value. Conversely, if a valid $x$ exists,
then $-(2x + 17)$ is in $A$ by the validity of $x$ and in $B$ by definition.

The time complexity of this algorithm is $\mathcal O(n\log n)$.
Indeed, in the worst case where no valid value is found,
we iterate through binary search $n$ times, and each application
of binary search takes $\mathcal O(\log n)$ operations
because, by assumption, the binary search is being applied
to the $n$-element array $B$ in the worst case where it won't find the key.
Additionally, \textsc{FindX-Better} takes $\Theta(n)$ extra space
due to generating the $n$-element array $B$.

\begin{alg}{\textsc{FindX-Best($A[1\dots n]$)}}
	\ii valuefound $\gets$ false
	\ii left $\gets 1$
	\ii right $\gets n$
	\ii while right $>$ 0 and left $\le n$ and not valuefound do
	\ii \qquad if $A[\text{left}] = -(2\cdot A[\text{right}] + 17)$ then
	\ii \qquad \qquad print($A[\text{right}]$)
	\ii \qquad \qquad valuefound $\gets$ true
	\ii \qquad \qquad break
	\ii \qquad else if $A[\text{left}] > -(2\cdot A[\text{right}] + 17)$ then
	\ii \qquad \qquad right $\gets \text{right} - 1$
	\ii \qquad else left $\gets \text{left} + 1$
	\ii if not valuefound then
	\ii \qquad print(``no valid value found'')
\end{alg}

To argue that this algorithm works we show that
the algorithm will always print out the largest
valid value (if it exists) in $A$.
Clearly if no valid value exists
then eventually either right will equal $0$ or left will equal $n+1$
and the algorithm will print ``no valid value found.''

Now suppose that $x$ is a valid value whose index in $A$ is $a$
(if $x$ appears more than once in $A$, pick $m$ to be as large as possible),
and suppose that the index of the first occurrence of $-(2x + 17)$ in $A$
is $b$. Eventually, right will equal $a$; the algorithm can't halt
before this point because there are no indices larger than $a$
corresponding to valid values. Then, left cannot be larger than $b$.
To see this, suppose left were larger than $b$.
left must have then previously equaled $b$, at which point
right would be greater than $a$ (if it equaled $a$ then
$A[a]$ would be detected as a valid value).
Observe the inequality
\[A[b] = -(2\cdot A[a] + 17) > -(2\cdot A[c] + 17)\]
for all $c > a$, since $A[c]$ would be strictly greater than $A[a]$
(recall that $a$ is the largest index giving value $A[a]$).
Thus, when $\text{left} = b$ and $\text{right} > a$,
the algorithm would repeatedly decrease right by $1$ until it equals $a$
rather than increasing left.
Now we've concluded that when right equals $a$, left is at most $b$.
Symmetric to the above logic, if left is equal to $b$ then we're done;
and, for every $c < b$,
\[A[c] < A[b] = -(2\cdot A[a] + 17)\]
(recalling $b$ is the smallest index giving value $A[b]$),
so when left is smaller than $b$ and right equals $a$,
the algorithm keeps incrementing left until left equals $b$
at which point $A[a]$ is found as a valid value.

The while loop takes at most $2n-1$ steps to terminate,
since right is allowed to be decremented at most $n-1$ times and
left is allowed to be incremented at most $n-1$ times.
The body of the while loop runs in constant time,
which means that this algorithm is $\mathcal O(n)$.
Furthermore, the only extra space is taken by the variables
valuefound, left, and right, amounting to $\Theta(1)$ extra space.

\pagebreak

\section{Problem 2}

If $A = [1,2,3,4,5]$,
then these are the states the array passes through.
Here, $i$ is increasing top to bottom
and $j$ from left to right; the cell with index $i$ is colored red
while the cell with index $j$ is colored green (except when $i = j$,
which gets blue).
\begin{center}
	\begin{asy}
		unitsize(0.5cm);
		int[] A = {1, 2, 3, 4, 5};
		pair s = (0,0);
		draw((s+(0,0))--(s+(5,0))--(s+(5,1))--(s+(0,1))--cycle);
		for (int i = 1; i < 5; ++i) draw((s+(i,0))--(s+(i,1)));
		for (int i = 0; i < 5; ++i) label(string(A[i]), s+(i+0.5,0.5));
		label("initial array", s + (2.5,-0.5));
		
		for (int i = 0; i < 5; ++i) {
			for (int j = 0; j < 5; ++j) {
				pair s = (6*(j-2),-2.4*(i+1));
				if (i == j) {
					fill((s+(i,0))--(s+(i+1,0))--(s+(i+1,1))--(s+(i,1))--cycle,blue+opacity(0.3));
				}
				else {
					fill((s+(i,0))--(s+(i+1,0))--(s+(i+1,1))--(s+(i,1))--cycle,red+opacity(0.3));
					fill((s+(j,0))--(s+(j+1,0))--(s+(j+1,1))--(s+(j,1))--cycle,heavygreen+opacity(0.3));
				}
				if (A[i] < A[j]) { // swap
					int tmp = A[i];
					A[i] = A[j];
					A[j] = tmp;
				}
				draw((s+(0,0))--(s+(5,0))--(s+(5,1))--(s+(0,1))--cycle);
				for (int i = 1; i < 5; ++i) draw((s+(i,0))--(s+(i,1)));
				for (int i = 0; i < 5; ++i) label(string(A[i]), s+(i+0.5,0.5));
				label("("+string(i+1)+","+string(j+1)+")", s + (2.5,-0.5));
			}
		}
	\end{asy}
\end{center}

If $A = [1,2,4,3,2]$:
\begin{center}
	\begin{asy}
		unitsize(0.5cm);
		int[] A = {1, 2, 4, 3, 2};
		pair s = (0,0);
		draw((s+(0,0))--(s+(5,0))--(s+(5,1))--(s+(0,1))--cycle);
		for (int i = 1; i < 5; ++i) draw((s+(i,0))--(s+(i,1)));
		for (int i = 0; i < 5; ++i) label(string(A[i]), s+(i+0.5,0.5));
		label("initial array", s + (2.5,-0.5));
				
		for (int i = 0; i < 5; ++i) {
			for (int j = 0; j < 5; ++j) {
				if (A[i] < A[j]) { // swap
					int tmp = A[i];
					A[i] = A[j];
					A[j] = tmp;
				}
				pair s = (6*(j-2),-2.4*(i+1));
				if (i == j) {
				fill((s+(i,0))--(s+(i+1,0))--(s+(i+1,1))--(s+(i,1))--cycle,blue+opacity(0.3));
				}
				else {
				fill((s+(i,0))--(s+(i+1,0))--(s+(i+1,1))--(s+(i,1))--cycle,red+opacity(0.3));
				fill((s+(j,0))--(s+(j+1,0))--(s+(j+1,1))--(s+(j,1))--cycle,heavygreen+opacity(0.3));
				}
				draw((s+(0,0))--(s+(5,0))--(s+(5,1))--(s+(0,1))--cycle);
				for (int i = 1; i < 5; ++i) draw((s+(i,0))--(s+(i,1)));
				for (int i = 0; i < 5; ++i) label(string(A[i]), s+(i+0.5,0.5));
				label("("+string(i+1)+","+string(j+1)+")", s + (2.5,-0.5));
			}
		}
	\end{asy}
\end{center}

This algorithm always results in a sorted array.
To prove this, we will prove the stronger statement that
the $k$th iteration of the outer loop (i.e. after the pairs
$(k, 1)$, $(k, 2)$, \dots, $(k, n)$ are all processed),
the subarray $A[1\dots k]$ is sorted and
$A[k]$ is the maximal value in the array.

We shall prove this by induction, starting with the base case of $k=1$.
In this case, whenever $j$ reaches the first index for which $A[j]$ is maximal,
$A[1] < A[j]$ must hold, so the values at indices $1$ and $j$ swap;
after this point there are no further swaps because
there are no larger values than $A[j]$ in the array.
Of course, the $1$-element array $A[1\dots 1]$ is trivially sorted.

Now suppose that $A[1\dots k]$ is sorted and $A[k] = n$,
and we are on the $(k+1)$th iteration of the loop (i.e. $i = k+1$).
``Weaken'' the swapping condition to say $A[i]$ and $A[j]$ are swapped
if $A[i]< A[j]$ or $A[i]=A[j]$ and $i>j$;
obviously swapping the identical numbers doesn't affect anything,
so this is just a cosmetic change.
Then, let $j\in\{1, \dots, k\}$ be the first index for which $A[k+1] \le A[j]$.
The first swap performed is between $A[k+1]$ and $A[j]$.
If $j\le k$, the new value $A[j]$ at index $k+1$ is now at most $A[j+1]$,
so the next operation will be swapping the new $A[k+1]$ (which is $A[j]$)
and $A[j+1]$; then the new new $A[k+1]$ (equals to the initial $A[j+1]$)
swaps with $A[j+2]$; and this will keep going until reaching $k+1$.
The sequence of swaps thus results in the following situation:
\begin{center}
	\begin{tabular}{c|c|c|c|c|c}
		$j$ & $j+1$ & $j+2$ & \dots & $k$ & $k+1$\\
		\hline
		$A[k+1]$ & $A[j]$ & $A[j+1]$ & \dots & $A[k-1]$ & $A[k]$
	\end{tabular}
\end{center}
Here the array values $A[t]$ are the values of the array
prior to all of the swaps listed out above.
After this point no other swaps are made because $A[k]$ is maximal
and the index is larger than $k+1$.
By assumption $A[1] \le A[2]\le \cdots \le A[k+1]\le A[j]\le\cdots\le A[k]$
so the new $A[1\dots (k+1)]$ is sorted and the new $A[k+1]$
(equal to $A[k]$ in the old array) is maximal.

\section{Problem 3}
\begin{alg}{\textsc{CountInversions-NonRecursive($A[1\dots n]$)}}
	\ii count $\gets 0$
	\ii for $i\gets 1$ to $n$ do
	\ii \qquad for $j\gets i+1$ to $n$ do
	\ii \qquad\qquad if $A[i] > A[j]$ then count $\gets \text{count} + 1$
	\ii return count 
\end{alg}

This algorithm iterates through all pairs of indices
$(i, j)$ with $1 \le i<j\le n$ and checks whether each pair is an inversion,
incrementing the count whenever it finds one.

The algorithm has a time complexity of $\Theta(n^2)$ since it always
processes $\binom n2 = \frac{n^2 - n}{2}$ pairs of indices,
and it uses $\Theta(1)$ extra space for the count variable.

\begin{alg}{\textsc{CountInversions-Recursive($A[1\dots n]$)}}
	\ii if $n = 1$ then
	\ii \qquad return $0$
	\ii else
	\ii \qquad terminalcontrib $\gets 0$
	\ii \qquad for $i\gets 1$ to $n-1$ do
	\ii \qquad\qquad if $A[i] > A[n]$ then terminalcontrib
	$\gets\text{terminalcontrib} +1$
	\ii \qquad return \textsc{CountInversions-Recursive($A[1\dots (n-1)]$)}
	+ terminalcontrib
\end{alg}

This recursive algorithm computes the number of inversions
by counting the number of inversions among the first $n-1$ elements of $A$
and adding this to the number of inversions formed by the terminal element;
for the base case of $n=1$, there are of course trivially zero inversions.

The time complexity is
\[T(n) = \begin{cases}
	\Theta(1) & n=1\\
	T(n-1) + \Theta(n) & n>1
\end{cases}\]
which amounts to $T(n) \in \Theta(n^2)$,
where the $+\Theta(n)$ comes from the for loop that computes terminalcontrib.
The space complexity is
\[S(n)=\begin{cases}
	\mathcal O(1) & n=1\\
	S(n-1) + \Theta(1) & n>1
\end{cases}\]
with the $+\Theta(1)$ coming from keeping track of terminalcontrib,
which means $S(n)\in\Theta(n)$.

\pagebreak
\section{Problem 4}
\begin{alg}{\textsc{PrefixSum-NonRecursive($A[1\dots n]$)}}
	\ii create array $P[1\dots n]$
	\ii $P[1]\gets A[1]$
	\ii for $i\gets 2$ to $n$ do
	\ii \qquad $P[i]\gets P[i-1]+A[i]$
\end{alg}

This algorithm computes the prefix sums iteratively
by initializing $P[1] = A[1]$ and using the identity
\[P[i-1] + A[i] = (A[1] + A[2] + \cdots + A[i-1]) + A[i] = P[i].\]

The time complexity is $\Theta (n)$, since the loop runs $n$ times
and each run of the loop takes constant time.
The space complexity is $\Theta(n)$, coming from the space taken up by
the prefix array $P[1\dots n]$.


\begin{alg}{\textsc{PrefixSum-Recursive($A[1\dots n]$)}}
	\ii if $n = 1$ then
	\ii \qquad return $A[1\dots 1]$
	\ii else
	\ii \qquad create array $P[1\dots n]$
	\ii \qquad $P[1\dots(n-1)] \gets
	\textsc{PrefixSum-Recursive($A[1\dots (n-1)]$)}$
	\ii \qquad $P[n]\gets P[n-1] + A[n]$
	\ii \qquad return $P[1\dots n]$
\end{alg}

This algorithm computes the prefix sums by first calling
\textsc{PrefixSum-Recursive} to compute the prefix sums for
the first $n-1$ elements of the array, and then computes
the $n$th prefix sum with the aforementioned identity $P[i-1] + A[i] = P[i]$.

The time complexity is
\[T(n) =\begin{cases}
	T(1) & n=1\\
	T(n-1) + \Theta(1) & n > 1
\end{cases}\]
which implies that $T(n)\in \Theta(n)$.
The space complexity is
\[S(n)=\begin{cases}
	S(1) & n=1\\
	S(n-1) + \Theta(n) & n>1
\end{cases}\]
where the $\Theta(n)$ comes from creating a new array
every time \textsc{PrefixSum-Recursive} is called,
so the overall space complexity is $S(n)\in\Theta(n^2)$.


\pagebreak
\section{Problem 5}
\begin{alg}{\textsc{IsValidString($S[1\dots 2n]$)}}
	\ii // returns true if $S$ is valid and false otherwise
	\ii isvalid $\gets$ true
	\ii lcount $\gets 0$
	\ii rcount $\gets 0$
	\ii for $i\gets 1$ to $2n$ do
	\ii \qquad if $S[i] = \text{`('}$ then lcount $\gets\text{lcount} + 1$
	\ii \qquad if $S[i] = \text{`)'}$ then rcount $\gets\text{rcount} + 1$
	\ii \qquad if lcount $<$ rcount then isvalid $\gets$ false
	\ii return isvalid
\end{alg}

The main idea of this algorithm is that a string is valid
if and only if, as one reads the string from left to right,
the number of left parentheses (lcount) is always at least as large as
the number of right parentheses (rcount). Clearly this is a necessary condition;
to see that it is sufficient, one can use induction on $n$:
the base case $n=1$ is trivial (the only valid string is ()),
and in general if one deletes a matching pair of adjacent parentheses
then one gets a string of length $2(n-1)$ satisfying the same condition,
which by hypothesis is valid and remains valid
upon reintroduction of the deleted pair.

The time complexity here is $\Theta(n)$, as the loop runs exactly $2n$ times
and the body of the loop runs in constant time.
The space complexity is $\Theta(1)$ because the only extra space used
is to store the variables isvalid, lcount, and rcount.

\section{Problem 6}
\begin{alg}{\textsc{GCD-NonRecursive($a,b$)}}
	\ii $g\gets 1$
	\ii for $i\gets 1$ to $\min(a, b)$ do
	\ii \qquad if $a\%i = 0$ and $b\%i = 0$ then
	\ii \qquad\qquad $g\gets i$
	\ii return $g$
\end{alg}

This algorithm just manually checks all numbers
not larger than both $a$ and $b$ for whether they are common divisors
of $a$ and $b$.

The time complexity is $\Theta(\min(a,b))$ because the for loop
always runs $\min(a,b)$ times (and the body of the for loop
runs in constant time). The space complexity is $\Theta(1)$,
with the only extra space used being $g$.

\begin{alg}{\textsc{GCD-Recursive($a,b$)}}
	\ii for $i\gets 2$ to $\min(a, b)$ do
	\ii \qquad if $a\%i = 0$ and $b\%i = 0$ then
	\ii \qquad\qquad return $i\cdot\textsc{GCD-Recursive}(a/i,b/i)$
	\ii return 1
\end{alg}

This algorithm is a twist on the non-recursive edition:
it checks for a divisor larger than $1$ and, upon finding such a divisor $d$,
computes $d\cdot \gcd(a/d, b/d)$. However, this is not much of an improvement
to the time complexity, since the case where $\gcd(a, b) = 1$
still requires $\min(a, b)$ operations (and it's easy to see that
the time complexity is bounded above by $\mathcal O(\min(a,b))$),
so the algorithm is still $\mathcal O(\min(a,b))$.

In terms of space complexity the algorithm has a different worst case,
namely when the number of recursive calls made is as large as possible.
This necessitates that $a$ and $b$ share a lot of small prime factors,
so the worst case is when $b = 2^n$ and $b$ divides $a$,
in which case the extra space is on the order of $\log n$, so in general
the space complexity of the algorithm is $\mathcal O(\log n)$.

\pagebreak
\section{Problem 7}
\begin{alg}{\textsc{EvaluatePoly-Naive($A[0\dots n],x$)}}
	\ii value $\gets 0$
	\ii for $i\gets 0$ to $n$ do
	\ii \qquad powofx $\gets1$
	\ii \qquad for $j\gets 1$ to $i$ do
	\ii \qquad \qquad powofx $\gets\text{powofx}\cdot x$
	\ii \qquad value $\gets\text{value} + \text{powofx}\cdot A[i]$
	\ii return value
\end{alg}

This algorithm just computes each term of the polynomial
by calculating $x^i$ in $\mathcal O(i)$ fashion and then
adds everything together.
The time complexity is $\mathcal O(n^2)$ because of the nested for loops,
and the space complexity is just $\Theta(1)$.

\begin{alg}{\textsc{EvaluatePoly-Better($A[0\dots n],x$)}}
	\ii value $\gets 0$
	\ii for $i\gets 0$ to $n$ do
	\ii \qquad powofx $\gets1$
	\ii \qquad exp $\gets i$
	\ii \qquad $y\gets x$
	\ii \qquad while $\text{exp}>0$ do
	\ii \qquad\qquad if $\text{exp}\%2=1$ then
	$\text{powofx}\gets\text{powofx}\cdot y$
	\ii \qquad\qquad $y \gets y\cdot y$
	\ii \qquad\qquad exp $\gets\text{exp}/2$
	\ii \qquad value $\gets\text{value} + \text{powofx}\cdot A[i]$
	\ii return value
\end{alg}

This algorithm runs almost identically to the previous one,
but with an improvement on the computation of $x^i$
by using binary exponentiation instead of linear exponentiation.
The binary exponentiation improves the complexity to
$\mathcal O(n\log n)$, and the space complexity remains as $\Theta(1)$.

\begin{alg}{\textsc{EvaluatePoly-Best($A[0\dots n],x$)}}
	\ii value $\gets 0$
	\ii powofx $\gets 1$
	\ii for $i\gets 0$ to $n$ do
	\ii \qquad value $\gets\text{value} + \text{powofx}\cdot A[i]$
	\ii \qquad powofx $\gets\text{powofx}\cdot x$
	\ii return value
\end{alg}

This algorithm directly improves on the naive one
by not throwing away power of $x$ calculated in the preceding iteration,
so that in a given run of the loop $x^i$ can be computed as
$x\cdot x^{i-1}$ where we've held on to the value of $x^{i-1}$ from earlier.
This makes the algorithm $\mathcal O(n)$ now, since the body of the for loop
now runs in constant time.
The space complexity is still $\Theta(1)$, like the
two previous iterations of this algorithm.

\pagebreak
\section{Problem 8}
\begin{alg}{\textsc{PairSum-Naive($A[1\dots n],k$)}}
	\ii // returns $(i, j)$ if $i<j$ and $A[i] + A[j] = k$,
	or $(-1, -1)$ if no such $(i, j)$ exists
	\ii for $i\gets 1$ to $n$ do
	\ii \qquad for $j\gets i+1$ to $n$ do
	\ii \qquad\qquad if $A[i] + A[j] = k$ return $(i, j)$
	\ii return $(-1,-1)$
\end{alg}

This is a brute force algorithm that checks each pair $(i, j)$
with $1\le i<j\le n$ to see if $A[i] + A[j] = n$;
there are $\binom n2\in\mathcal O(n^2)$ such pairs
and each pair is processed in constant time, so the time complexity is
$\mathcal O(n^2)$. The space complexity is $\Theta(1)$;
nothing interesting happening here.

\begin{alg}{\textsc{PairSum-Better($A[1\dots n],k$)}}
	\ii // returns $(i, j)$ if $i<j$ and $A[i] + A[j] = k$,
	or $(-1, -1)$ if no such $(i, j)$ exists
	\ii create array $B[1\dots n]$
	\ii for $i\gets 1$ to $n$ do
	\ii \qquad $B[n+1-i] = k - A[i]$
	\ii for $i\gets 1$ to $n$ do
	\ii \qquad $j\gets n+1-\textsc{BinarySearch}(B[1\dots (n-i)],A[i])$
	\ii \qquad if $j > 0$ then return $(i, j)$
	\ii return $(-1, -1)$
\end{alg}

The idea here is that for each $A[i]$, we are searching among
all the elements appearing after $A[i]$ to look for $k-A[i]$;
the way this is done mechanically is via the array $B$,
which stores all the $k-A[i]$s in reverse order (reverse to ensure
that $B[1\dots n]$ is sorted instead of reverse-sorted).

The time complexity of this algorithm is $\mathcal O(n\log n)$
due to the binary search being called $n$ times;
this asymptotic bound is sharp in the worst case where
there is no $(i, j)$ pair with $i<j$ and $A[i]+A[j]=k$.
The space complexity is $\Theta(n)$ because the algorithm
requires the production of the auxiliary array $B$ with length $n$.

\begin{alg}{\textsc{PairSum-Best($A[1\dots n],k$)}}
	\ii // returns $(i, j)$ if $i<j$ and $A[i] + A[j] = k$,
	or $(-1, -1)$ if no such $(i, j)$ exists
	\ii $\ell\gets 1$
	\ii $r\gets n$
	\ii while $\ell < r$ do
	\ii \qquad if $A[\ell] + A[r] = k$ then return $(\ell, r)$
	\ii \qquad else if $A[\ell] + A[r] > k$ then $r\gets r-1$
	\ii \qquad else $\ell\gets \ell + 1$
	\ii return $(-1, -1)$
\end{alg}

This algorithm is in very similar spirit to the solution to Problem 1(iii).
The idea is that increasing $\ell$ will increase
the sum $A[\ell] + A[r]$, while decreasing $r$ will decrease the sum
because of the sorted nature of $A$. If an $(i,j)$ pair with
the desired properties doesn't exist, then it's clear that the algorithm
will work because eventually $\ell$ will exceed $r$ (after at most $n$ steps,
since $r-\ell$ strictly decreases after each iteration of the loop).

Otherwise, suppose that an $(i, j)$ pair exists,
and take $i$ to be the smallest $i$ appearing among all these pairs,
so that $j$ can be taken as the largest $j$ among all these pairs
(indeed, if there were a pair $(i', j')$ with $j' > j$ then either
$A[j'] = A[j]$, so we can just replace $j\gets j'$, or $A[j'] > A[j]$
in which case $A[i'] < A[i]$ implies $i' < i$, contradiction).
Then, by a similar argument to the one made in Problem 1(iii)
(here with $k-x$ instead of the $-(2x+17)$ expression),
we can argue that eventually either $\ell$ will equal $i$
or $r$ will equal $j$, and then in the first case
we must have $r$ repeatedly decrease until it equals $j$,
or in the second case we must have $\ell$ repeatedly incremented
until it equals $i$, so the algorithm stops
once it finds the valid pair $(\ell, r) = (i, j)$.

The time complexity is $\mathcal O(n)$, since the while loop runs
at most $n$ times, as remarked in the first paragraph,
and the body of the while loop runs in constant time.
The space complexity is $\Theta(1)$, as there are only some
extra auxiliary variables ($\ell$ and $r$).

\pagebreak
\section{Problem 9}

\begin{alg}{\textsc{CountFrequency-Naive($A[1\dots n],k$)}}
	\ii count $\gets 0$
	\ii for $i\gets 1$ to $n$ do
	\ii \quad if $A[i] = k$ then count $\gets\text{count} + 1$
	\ii return count
\end{alg}

This just checks every element of $A$ for whether it equals $k$
and increments the count when appropriate.
The time complexity is clearly $\mathcal O(n)$
and the space complexity clearly $\Theta(1)$.

\begin{alg}{\textsc{CountFrequency-Best($A[1\dots n], k$)}}
	\ii first $\gets 0$; last $\gets -1$
	\ii if $A[1] = k$ then first $\gets 1$
	\ii else
	\ii \qquad $\ell\gets 1$; $r\gets n$
	\ii \qquad while $\ell\le r$ do
	\ii \qquad\qquad $m\gets (\ell + r)/2$
	\ii \qquad\qquad if $A[m] = k$ and $A[m-1] \ne k$ then
	\ii \qquad\qquad\qquad first $\gets m$
	\ii \qquad\qquad\qquad break
	\ii \qquad\qquad else if $A[m] < k$ then $\ell\gets m+1$
	\ii \qquad\qquad else $r\gets m-1$
	\ii if $A[n] = k$ then last $\gets n$
	\ii else
	\ii \qquad $\ell\gets 1$; $r\gets n$
	\ii \qquad while $\ell\le r$ do
	\ii \qquad\qquad $m\gets (\ell + r)/2$
	\ii \qquad\qquad if $A[m] = k$ and $A[m+1] \ne k$ then
	\ii \qquad\qquad\qquad last $\gets m$
	\ii \qquad\qquad\qquad break
	\ii \qquad\qquad else if $A[m] > k$ then $r\gets m-1$
	\ii \qquad\qquad else $\ell\gets m+1$
	\ii return $\text{last} - \text{first} + 1$
\end{alg}

This algorithm finds the first instance of $k$ in the array
and the last instance of $k$ in the array, and then uses these indices
to compute the number of times $k$ appears (since every number
between the first and last $k$ is equal to $k$ because $A$ is sorted).

The time complexity of this algorithm is $\mathcal O(\log n)$
because in both while loops, $r - \ell$ gets halved
in each iteration of the loop, so the loops run $\mathcal O(\log n)$ times
with constant time each iteration.
The space complexity is obviously $\Theta(1)$.

\pagebreak
\section{Problem 10}

To implement stacks using queues, we will treat
the back of the queue as the top of the stack.
\begin{alg}{\textsc{PushWithQueue}}
	\ii // given queue $Q$ behaving like a stack, push $e$ to the stack
	\ii $Q$.enqueue($e$)
\end{alg}
Since the back of the queue is the top of the stack, adding to the top
means just enqueuing $e$. Constant time, constant space.

\setcounter{algo}{0}
\begin{alg}{\textsc{PopWithQueue}}
	\ii // given queue $Q$ behaving like a stack, pop the top of the stack
	\ii if $Q$.size() = 0 then return null
	\ii create queue $R$
	\ii for $i\gets 1$ to $Q.\text{size}() - 1$ do
	\ii \qquad $R$.enqueue($Q$.dequeue())
	\ii popped $\gets$ $Q$.dequeue()
	\ii $Q\gets R$
	\ii return popped
\end{alg}

Here, we need to remove the element at the back of $Q$ and return it.
To do this we must dequeue $Q$.size() times, the latter of which
will be the popped value that we need to return.
The remaining values are kept in a second queue $R$,
which, because of how queues work, has its elements
in the same order that they appeared in $Q$.
Therefore we can simply relabel $R$ as $Q$ and then return the popped value.
This is linear in both time and space complexity,
owing to the creation of the extra queue.\\

For stacks, we will treat the top of the stack as the back of the queue
and the bottom of the stack as the tail.
\begin{alg}{\textsc{EnqueueWithStack}}
	\ii // given stack $S$ behaving like a queue, enqueue $e$ to the queue
	\ii $S$.push($e$)
\end{alg}

Symmetric to pushing with the queue, in order to enqueue $e$
we just push $e$ to the top of the stack (which is the back of the queue).
Again, constant time and space.

\setcounter{algo}{1}
\begin{alg}{\textsc{DequeueWithStack}}
	\ii // given stack $S$ behaving like a queue, dequeue the front of the queue
	\ii if $S$.size() = 0 then return null
	\ii create stack $T$
	\ii for $i\gets 1$ to $S.\text{size}() - 1$ do
	\ii \qquad $T$.push($S$.pop())
	\ii dequeued $\gets$ $S$.pop()
	\ii for $i\gets 1$ to $T.\text{size}()$ do
	\ii \qquad $S$.push($T$.pop())
	\ii return dequeued
\end{alg}

Here we have to do similar shenanigans with making a copy of the stack
and popping every element to dequeue the one at the bottom.
However, in this case, the storage unit $T$, being a stack,
has its elements in the opposite order that they appeared in $S$
(so the top of $T$ was the second element from the bottom of $S$,
while the bottom of $T$ was the top of $S$).
Thus we need to move the elements of $T$ back into $S$ as well.
This again ends up having linear time and space complexities.












\end{document}